\chapter{Финальное заключение Тома IX и программа Тома X}
\label{ch:v9-final-conclusion-tome10-program}

\section{Какие задачи ставились}

Том IX был открыт после того, как Том VIII построил полный корреляционный
процесс, но оставил независимыми четыре физические компоненты: концевой сектор,
энергетический масштаб $E_*$, безразмерную связь $\chi$ и закон транспорта.
Поэтому центральная задача состояла не в построении ещё одной допустимой
модели, а в выводе одного общего ограниченного снизу родительского функционала,
который совместно выбирает все четыре компоненты без подстановки наблюдаемых данных.

Эта задача была разложена на шесть проверяемых требований:
\begin{enumerate}
\item построить общий носитель, содержащий расширение концевого сектора как внутренний
сектор, а не как вручную добавленное внешнее пространство;
\item вывести из одного действия положительный $E_*$ и связь $\chi$;
\item получить примитив переноса, совместимый с GNVW-границей;
\item совместно определить стационарное состояние и неподвижную алгебру полного процесса;
\item получить неотрицательный полный физический гессиан после факторизации по
калибровочным направлениям;
\item вывести хотя бы одно слепое безразмерное следствие, не использованное
при выборе родительского функционала.
\end{enumerate}

\section{Ответ Тома IX}

Том IX проверял, существует ли один ограниченный снизу динамический
родительский функционал, совместно выбирающий концевой сектор, физический
масштаб $E_*$, связь $\chi$, перенос, стационарный процесс, полный гессиан
и слепое следствие.

На допущенных расширениях построена полная условная модель:
\begin{equation}
 \mathbf n_{IX}^{\rm cond}=(1,1,1,1,1,1)^T,
 \qquad N_{IX}^{\rm cond}=6/6.
 \label{eq:v9-final-conditional-status}
\end{equation}
Однако строгий физический реестр равен
\begin{equation}
 \mathbf n_{IX}^{\rm phys}=(1,0,1,1,0,0)^T,
 \qquad N_{IX}^{\rm phys}=3/6.
 \label{eq:v9-final-physical-status}
\end{equation}
Точная разность
\begin{equation}
 \mathbf d_{IX}=(0,1,0,0,1,1)^T,
 \qquad
 \mathbf n_{IX}^{\rm phys}+\mathbf d_{IX}=\mathbf n_{IX}^{\rm cond},
 \qquad \rank\operatorname{diag}(\mathbf d_{IX})=3
 \label{eq:v9-final-status-decomposition}
\end{equation}
состоит из физического происхождения масштаба и связи, физического гессиана
и безусловного слепого следствия.

\section{Что удалось достичь}

В Том IX входят следующие устойчивые результаты:
\begin{enumerate}
\item построены общий четырёхкомпонентный носитель и ограниченное снизу
условное семейство родительских функционалов;
\item найдены минимальный модуль концевого сектора, общее пространство
конфигураций и квантово-марковский процесс рождения, переводящий три типа
концевого сектора в единый динамический язык;
\item построено двунаправленное примитивное KMS-завершение на $M_6(\mathbb C)$
с единственным точным стационарным состоянием;
\item шесть параметров KMS разделены на два общих масштаба и четыре
относительные формы, благодаря чему точно локализована недоопределённость;
\item построен минимальный инвариантный механизм выбора всех четырёх форм
посредством логарифма определителя;
\item после явного добавления новой аксиомы логарифма определителя получен
общий дополненный родительский функционал с гессианом полного ранга и
изолированным минимумом;
\item из дополненной модели выведены точные условные отношения, отклики и
нулевая взвешенная дисперсия;
\item построена исчерпывающая карта запретительных результатов для физической
меры логарифма определителя, опорного состояния и абсолютного масштаба.
\end{enumerate}

Критерий повторного открытия требовал двух независимых пакетов происхождения. Итог равен
\begin{equation}
 \mathbf a_{\rm reopen}^{\rm phys}=(0,0)^T,
 \qquad N_{\rm reopen}^{\rm phys}=0/2.
 \label{eq:v9-final-reopening-status}
\end{equation}
Носитель гауссовой ковариации и опорное состояние OU существуют как условные
архитектуры, но сохраняют соответственно $E_*/\mu$ и $\delta/\gamma$.
Восемь кандидатов масштаба дали $0/8$; общая нулевая масштабная мода не снята.

\section{Чего достичь не удалось}

Ни одна проверенная конструкция не вывела абсолютный энергетический масштаб
$E_*$ из внутренних данных. Температура KMS, разрыв часов, масштаб отсечения,
обратный радиус и гауссов опорный масштаб задают лишь отношения и сохраняют
общее масштабное преобразование. Наблюдаемая масса не может использоваться
как механизм выбора, поскольку тогда целевой ответ уже содержится во входных данных.

Не удалось также вывести физическое происхождение члена логарифма определителя.
Вспомогательные фермионы, завершение BRST/Штюкельберга, физическая фермионная
петля, связанный резервуар, функционал влияния Келдыша, спектральная плотность
резервуара и аномалия меры либо требуют нового постулата, либо дают сокращение
определителей, либо
оставляют дополнительные свободные данные.

Следовательно, не получены три пункта строгого контракта: совместный
физический механизм выбора $E_*,\chi$, физический гессиан без новой аксиомы
и безусловное слепое следствие. Условная гауссова архитектура не исправляет
этот итог: её ковариация зависит от свободного отношения $\delta/\gamma$,
а общий носитель --- от $E_*/\mu$.

\begin{theorem}[Финальный ответ Тома IX]
Один математически замкнутый дополненный родительский функционал построен и
даёт условный счёт $6/6$. Но из корпуса Томов I--VIII и допустимых расширений
не выведены пакет физического происхождения масштаба и связи и пакет
физического происхождения логарифма определителя. Поэтому строгий физический
счёт остаётся $3/6$, а физический четырёхкомпонентный родительский функционал
не получен.
\end{theorem}

\begin{nogocriterion}[Финальная заморозка Тома IX]
Нельзя повышать условную модель логарифма определителя до физической теории,
использовать наблюдаемую массу как $\mu$, полагать $\delta=\gamma$ по
нормировке или переименовывать масштаб отсечения, радиус, разрыв часов и
температуру KMS в абсолютную энергию. Физический статус Тома IX заморожен.
\end{nogocriterion}

\section{Задачи следующего тома}

Продолжение на уровне конечномерного носителя, архитектуры QMS или ещё одной
нормировки повторяло бы закрытые ветви. Единственный новый класс механизма
--- квантовая ренормгрупповая динамика, следовая аномалия и размерная
трансмутация. Радиативно порождаемая шкала требует вычислимого квантового
эффективного действия и бегущих связей, как в механизме Коулмена--Вайнберга
\cite{v9-final-coleman-weinberg1973}.

Программа Тома X состоит из шести обязательств:
\begin{enumerate}
\item построить общий квантово-полевой носитель ренормгрупповой динамики;
\item вывести ненулевую $\beta$-функцию или следовую аномалию;
\item получить ренормгрупповой инвариантный масштаб трансмутации;
\item построить типизированное вложение этой шкалы в родительский функционал
KMS с гауссовым сектором;
\item вывести физическую меру и опорное состояние, порождающие член
логарифма определителя;
\item получить независимое от схемы перенормировки слепое следствие.
\end{enumerate}

\section{Ожидаемый результат Тома X}

Минимальное положительное ожидание состоит не в подгонке численного значения
$E_*$, а в доказательстве существования внутреннего механизма размерной
трансмутации. Затравочные носитель и действие должны быть заданы до вычисления
целевого значения; из них требуется получить ненулевой квантовый ход
ренормгруппы, её инвариантную шкалу и типизированную карту этой шкалы в уже
построенный гауссов сектор KMS.

Полный успех Тома X означает одновременное выполнение всех шести требований:
физический масштаб и связь выбираются квантовым эффективным родительским
функционалом, член логарифма определителя
возникает из его меры или интегрирования физических степеней свободы, полный
гессиан остаётся допустимым, а хотя бы одно независимое от схемы перенормировки безразмерное
следствие получается после, а не до выбора модели. Частичный успех допустим
как новый условный мост, но не должен повышаться до физического замыкания.

Отрицательное ожидание также содержательно: если вычисленная $\beta$-функция
нулевая, масштаб трансмутации остаётся зависимым от схемы или граничных
условий либо его типизированное вложение требует наблюдаемой массы, общее
нулевое масштабное направление сохраняется и программа физического
родительского функционала закрывается уже на квантово-полевом уровне.

Их матрица зависимостей фиксируется как
\begin{equation}
 D_X=
 \begin{pmatrix}
 1&0&0&0&0&0\\
 1&1&0&0&0&0\\
 0&1&1&0&0&0\\
 0&0&1&1&0&0\\
 0&0&0&1&1&0\\
 0&0&0&0&1&1
 \end{pmatrix},
 \qquad \rank D_X=6,
 \qquad \det D_X=1.
 \label{eq:v9-final-tome10-dependency}
\end{equation}
Требования полностью специфицированы, но ни одно ещё не построено:
\begin{equation}
 \mathbf n_X^{\rm specification}=(1,1,1,1,1,1)^T,
 \qquad
 \mathbf n_X^{\rm construction}=(0,0,0,0,0,0)^T.
 \label{eq:v9-final-tome10-status}
\end{equation}
Поэтому допуск Тома X означает исследовательскую программу, а не новую готовую
теорию.

\begin{nogocriterion}[Стоп-критерий Тома X]
Первый гейт Тома X обязан построить носитель, на котором квантовые поправки
и $\beta$-функция вычисляются из уже заданных представлений и действия.
Если носитель или затравочные связи выбираются после чтения целевого
масштаба, ветвь закрывается без дальнейшей численной подгонки.
\end{nogocriterion}

\begin{protocol}[Финальный статус Тома IX]\begin{enumerate}
\item условный Том IX: \textbf{$6/6$};
\item строгий физический Том IX: \textbf{$3/6$};
\item пакеты физического повторного открытия: \textbf{$0/2$};
\item физический четырёхкомпонентный родительский функционал: \textbf{не выведен};
\item Том IX: \textbf{завершён и физически заморожен};
\item спецификация Тома X: \textbf{$6/6$};
\item построение Тома X: \textbf{$0/6$};
\item следующий гейт: \textbf{допуск общего квантово-полевого носителя
ренормгрупповой динамики}.
\end{enumerate}\end{protocol}

Машинный сертификат сохранён в
\path{s2t_v9_final_conclusion_and_tome10_program_gate_results.json}.

\begin{thebibliography}{9}
\bibitem{v9-final-coleman-weinberg1973}
S.~Coleman, E.~Weinberg,
\emph{Radiative Corrections as the Origin of Spontaneous Symmetry Breaking},
Physical Review D 7 (1973), 1888--1910.
\end{thebibliography}